%% kapa: exam=1 round=28 year=2017 period=1 form=A subject=civil_law subject_ko=민법 seq=5 number=5
%% source: Q-Net 감정평가사 기출문제 (cst00309) artlSeq=5199603 file=extracted/1008/5199603/2017년 제28회 감정평가사 1차시험 문제지/1교시 A형.pdf
%% conversion: kordoc 4.24.51 -> tools/gen1.py (auto); review=proofread-sonnet; flags=none
\begin{question}{5}
민법상 법인에 관한 설명으로 옳지 않은 것은? (다툼이 있으면 판례에 따름)
\begin{choices}
  \item 사단법인의 정관에 다른 규정이 없는 한, 그 정관은 총사원 3분의 2 이상의 동의가 있는 때에 한하여 이를 변경할 수 있다.
  \item 법인은 법률의 규정에 의함이 아니면 성립하지 못한다.
  \item 법인의 목적 범위 외의 행위로 인하여 타인에게 손해를 가한 때에는 그 사항의 의결에 찬성하거나 그 의결을 집행한 사원, 이사 및 기타 대표자가 연대하여 배상하여야 한다.
  \item 사단법인의 사원의 지위는 양도 또는 상속할 수 없고, 이는 정관으로 달리 정할 수 없다.
  \item 이사가 수인(數人)인 경우에 법인의 사무집행은 정관에 다른 규정이 없는 한 이사의 과반수로써 결정한다.
\end{choices}
\end{question}
\answer{4}
